\section{FactMeta.sty}\label{sec:ref-factmeta}

Greppable \emph{fact} metadata -- ProblemMeta's little sibling, and like it
\emph{store-only} and nice-to-have. A fact is a bare snippet (a definition, a
formula) with no structure of its own; FactMeta lets you find one by what it
says without running \LaTeX. One command:

\begin{manexsource}
\FactMeta{
  source = {apex},
  kind   = {theorem},
  name   = {Power Rule},
  number = {2.1},
  topics = {differentiation, power-rule},
}
\end{manexsource}

Keys: \env{source}, \env{kind}, \env{name}, \env{number}, \env{topics}.
A fact has \emph{no} \env{type} or \env{difficulty} (those
describe a \emph{problem}'s form and challenge), and nothing extrinsic like
points. \env{topics} lists every subject the fact touches (and echoes its
folder); \env{source}/\env{number} are a provenance origin and its locator.
There is no \env{desc}: it was dropped 2026-08-13 when ProblemMeta's
\env{desc} was renamed \env{name}, the two being the same idea under two
names -- and no fact had ever used it.

\env{kind} and \env{name} are what let a fact render as the result it
actually \emph{is}. The fact bank is a catch-all -- definitions, theorems,
axioms -- so \env{kind} names the CourseBoxes environment the import site
should wrap the snippet in and \env{name} the title that environment
displays: the block above imports as ``Theorem 2.1 (Power Rule)''. Both are
optional; with no \env{kind} a fact imports bare, exactly as before.

\textbf{\env{kind} differs in nature from every other key here}, which is
worth knowing. The rest are descriptive -- search keys. \env{kind} is
\emph{functional}: it selects an environment, so a typo is not a fragmented
search but a fact that will not box. Two things contain that:
\cs{importfact} guards the environment at use and degrades to a bare import
with a warning rather than erroring, and \texttt{find-facts --vocab} lists
the kinds in the bank so drift (\texttt{thm} vs \texttt{theorem}) is visible
\emph{without} running \LaTeX. The vocabulary is closed in practice: a kind
is only usable if the course's text layer declared it with
\cs{NewTextResult}.

On the \LaTeX{} side the block is
\emph{validated but discarded}: a mistyped key errors at compile time, nothing is
typeset. The real reader is \file{find-facts}, which parses the block textually
off disk -- hence the same two disciplines (one key per line with braced values,
and \texttt{\#} written \texttt{\char`\\\#}) as ProblemMeta.

Design notes worth keeping in mind:

\begin{itemize}
  \item \textbf{Intrinsic only.} Source, number and topics belong to a fact
    forever. Anything set at the import site is not a key here.
  \item \textbf{Keys reset per block}, so a file that omits \env{number} never
    inherits a previous import's value when several facts land in one document.
  \item \textbf{Removable.} If tagging is abandoned, delete the package:
    Exercises' \cs{AtBeginDocument} stub swallows leftover \cs{FactMeta} blocks
    and every document still builds.
  \item \textbf{Import wrapper.} A picked fact is inserted as
    \cs{importfact}\texttt{[kind, name, number]\{DIR\}\{FILE\}} (defined
    in \file{Exercises.sty}, beside \cs{importproblem}). All three are captured
    into \cs{CurrentFactKind}, \cs{CurrentFactName} and \cs{CurrentFactNumber},
    and drive the wrapping box: ``Theorem 2.1 (Power Rule)''. With no \env{kind}
    there is no box, and a document can still print the name and number by
    redefining the \cs{ShowFactMeta} hook -- e.g.\ a formula sheet listing
    ``Power Rule 2.1'' beside each entry.
  \item \textbf{You do not retype these.} Ask \texttt{find-facts} for an
    import line with \texttt{--with kind,name,number} and it reads the block
    off disk, writing them into the emitted \cs{importfact} line. That the
    values land at the \emph{call site}
    is also what lets one import override them: the same fact can be a
    \env{theorem} in one course and a \env{keyidea} in another without editing
    the file.
  \item \textbf{An unnumbered fact still works.} With \env{name} but no
    \env{number} the import reaches for the kind's starred form, heading
    ``Theorem (Product, Quotient, and Chain Rules)''. A fact that resists a
    concise name can equally give \env{number} and no \env{name}.
\end{itemize}
